\documentclass[11pt,a4paper]{article}

\usepackage[T1]{fontenc}
\usepackage{newtxtext,newtxmath}
\usepackage[final]{microtype}
\usepackage[margin=1in]{geometry}
\let\openbox\relax
\usepackage{amsmath,amsthm,mathtools}
\usepackage{booktabs,array}
\usepackage[font=small,labelfont=bf,labelsep=period]{caption}
\usepackage[authoryear,round]{natbib}
\usepackage[colorlinks=true,allcolors=blue]{hyperref}
\usepackage[capitalize,nameinlink]{cleveref}
\usepackage{xurl}
\usepackage{setspace}

\onehalfspacing
\setlength{\parindent}{1.5em}
\setlength{\parskip}{0pt}
\emergencystretch=2em

\newtheorem{theorem}{Theorem}
\newtheorem{proposition}[theorem]{Proposition}
\newtheorem{corollary}[theorem]{Corollary}
\theoremstyle{definition}
\newtheorem{definition}[theorem]{Definition}
\theoremstyle{remark}
\newtheorem{remark}[theorem]{Remark}

\newcommand{\R}{\mathbb R}
\newcommand{\B}{\mathcal B}
\newcommand{\pos}[1]{[#1]_+}

\title{Certifying $\Gamma$-Robust Discrete Pricing via Full-Breakpoint\\
Multiple-Choice Knapsack Decomposition\\[0.7em]
\large V3 Legacy Reconstruction}
\author{Zi Yuan Eric Shao\\
\small Department of Mathematics, ETH Z\"urich}
\date{Original V3 research line: May 2026\\
Legacy reconstruction compiled: July 2026}

\hypersetup{
  pdftitle={Certifying Gamma-Robust Discrete Pricing via Full-Breakpoint Multiple-Choice Knapsack Decomposition - V3 Legacy Reconstruction},
  pdfauthor={Zi Yuan Eric Shao},
  pdfsubject={Transparent reconstruction from the V2 arXiv source and final V3 code archive}
}

\begin{document}
\maketitle

\begin{center}
\fbox{\begin{minipage}{0.91\textwidth}
\textbf{Archival status.}
The original V3 manuscript source and compiled PDF were stored only in an
ignored local artifact directory and are no longer available. This document
is a transparent legacy reconstruction from the exact arXiv v1 source of the
preceding paper, the final V3 code snapshot at Git commit
\texttt{2e2829e55243b846e5944965a839eb9b28dd5e9c}, and that snapshot's
\texttt{README}, \texttt{REVISION\_HISTORY}, \texttt{SUBMISSION}, tests, and
experiment drivers. It is not represented as a byte-identical copy of the
lost manuscript. Statements unsupported by preserved outputs are identified
as scope limits rather than recreated numerical claims.
\end{minipage}}
\end{center}

\begin{abstract}
We reconstruct the final scientific position of the V3 robust discrete-pricing
research line. A finite price menu with a portfolio margin requirement reduces
exactly to a multiple-choice knapsack problem. Under an integer
cardinality-budget uncertainty set, the coupled robust margin constraint is
equivalent to a finite union of fixed-threshold MCKPs over the complete set of
original deviation breakpoints. For every threshold, upper-hull geometry gives
an exact LP bound and a one-item additive rounding certificate. A globally
certifying branch-and-bound must preserve original integer options, initialize
a valid upper-bound record for every threshold, and validate every incumbent
against the original robust constraint. The V3 implementation further provides
an exact increasing-threshold sweep that updates transformed coefficients while
reusing a hull only when its point set is unchanged. This reconstruction states
the corresponding validity invariants, separates exact procedures from
diagnostic reductions, and records the limits of the surviving empirical
evidence. Its purpose is archival continuity between the public March 2026
paper and the independent V4 group-envelope pivot.
\end{abstract}

\noindent\textbf{Keywords:} robust optimization; discrete pricing;
multiple-choice knapsack; exact decomposition; parametric sweep

\section{Research position and provenance}

The V3 line studies a discrete portfolio-pricing model in which one price is
selected from each product menu. Portfolio revenue is maximized subject to a
minimum margin ratio and price-admissibility restrictions. The March 2026
paper established the exact reduction to a multiple-choice knapsack problem
(MCKP), the fixed-threshold LP geometry, and an additive one-item rounding
bound. V3 shifted the center of gravity from a fast heuristic narrative to a
certifying decomposition framework:

\begin{enumerate}
\item enumerate the complete original deviation-breakpoint set;
\item construct a valid fixed-threshold LP upper bound at every breakpoint;
\item maintain globally valid lower and upper bounds under time or node limits;
\item preserve below-hull integer options in exact search;
\item update the fixed-threshold data in an exact parametric sweep; and
\item treat reduced threshold sets and local-budget extensions as separately
scoped diagnostics unless their own exactness conditions hold.
\end{enumerate}

This is not a claim that threshold decomposition, MCKP hull geometry, or
Lagrangian relaxation is new. The contribution of the V3 program is the
combination of the pricing reduction, explicit robust and rounding
certificates, and defensible global-gap accounting in a reproducible solver.

\section{Discrete pricing and MCKP reduction}

Let $i\in\{1,\ldots,n\}$ index products and $J_i$ the admissible finite menu of
product $i$. Option $(i,j)$ has price $p_{ij}>0$, exposure $w_i\ge0$, reference
price $\bar p_i>0$, and nominal demand $\hat g_{ij}\ge0$. Binary variable
$x_{ij}$ selects one option:
\[
  \sum_{j\in J_i}x_{ij}=1,\qquad x_{ij}\in\{0,1\}.
\]
Define nominal revenue and reference-denominator coefficients
\[
  v_{ij}:=w_i p_{ij}\hat g_{ij},\qquad
  q_{ij}:=w_i\bar p_i\hat g_{ij}.
\]
For a required margin ratio $\Delta>0$, nominal feasibility is
\[
  \sum_{i,j}(v_{ij}-\Delta q_{ij})x_{ij}\ge0.
\]
Writing $s_{ij}:=v_{ij}-\Delta q_{ij}$ and
$s_i^*:=\max_j s_{ij}$ gives nonnegative transformed costs
$c_{ij}:=s_i^*-s_{ij}$ and capacity $C:=\sum_i s_i^*$. Thus the nominal model
is exactly
\begin{align*}
  \max_x\quad &\sum_{i,j}v_{ij}x_{ij}\\
  \text{s.t.}\quad &\sum_{i,j}c_{ij}x_{ij}\le C,\\
  &\sum_jx_{ij}=1\quad(i=1,\ldots,n),\qquad x\in\{0,1\}.
\end{align*}

\begin{proposition}[Exact nominal reduction]
The baseline transformation preserves every feasible price selection and its
objective value. Hence the finite-menu pricing model with a positive
denominator and a single portfolio ratio floor is an MCKP.
\end{proposition}

\begin{proof}
For every exactly-one selection,
$\sum_{i,j}s_{ij}x_{ij}=\sum_i s_i^*-\sum_{i,j}c_{ij}x_{ij}$.
The margin residual is nonnegative if and only if the transformed cost is at
most $C$. The objective and group equations are unchanged.
\end{proof}

\section{Budgeted uncertainty and full breakpoints}

Let $t_{ij}\in\R$ be the signed option-level margin exposure produced by an
adverse demand perturbation, and set $d_{ij}:=|t_{ij}|$. Under an integer
budget $\Gamma\in\{0,\ldots,n\}$, at most $\Gamma$ selected groups receive
their full adverse deviation. The robust margin residual is
\[
  \sum_{i,j}s_{ij}x_{ij}
  -\max_{\substack{0\le u_i\le1\\\sum_i u_i\le\Gamma}}
    \sum_i u_i\sum_jd_{ij}x_{ij}.
\]
LP duality for the inner budget problem yields
\[
 \max_{0\le u\le1,\ \mathbf 1^\top u\le\Gamma}\sum_i d_i u_i
 =\min_{\theta\ge0}
   \left\{\Gamma\theta+\sum_i\pos{d_i-\theta}\right\},
 \qquad d_i:=\sum_jd_{ij}x_{ij}.
\]
For fixed $\theta$, define
\[
  s_{ij}^{\theta}:=s_{ij}-\pos{d_{ij}-\theta}.
\]
The complete original candidate set is
\[
  \B:=\{0\}\cup\{d_{ij}:i=1,\ldots,n,\ j\in J_i\}.
\]

\begin{theorem}[Finite full-breakpoint decomposition]
A binary selection is robust feasible if and only if there exists
$\theta\in\B$ such that
\[
  \sum_{i,j}s_{ij}^{\theta}x_{ij}-\Gamma\theta\ge0.
\]
Consequently, the robust problem is the maximum over a finite family of
fixed-$\theta$ MCKPs.
\end{theorem}

\begin{proof}
For a fixed selection, the protection function
$h(\theta)=\Gamma\theta+\sum_i\pos{d_i-\theta}$ is convex and piecewise linear.
Its breakpoints are zero and the selected deviations. A minimizer therefore
exists at one of those values, all of which belong to the global set $\B$.
Substitution into the robust residual gives the stated equivalence.
\end{proof}

The word \emph{original} is essential. Removing an option before constructing
$\B$ can remove the only threshold supporting an optimal integer selection.
V3 therefore permits reduced threshold sets only as diagnostics or heuristics,
not as inputs to a global exactness claim.

\section{Fixed-threshold hull certificate}

For fixed $\theta\in\B$, let
\[
  s_i^*(\theta):=\max_j s_{ij}^{\theta},\qquad
  c_{ij}^{\theta}:=s_i^*(\theta)-s_{ij}^{\theta},\qquad
  C(\theta):=\sum_i s_i^*(\theta)-\Gamma\theta.
\]
If $C(\theta)<0$, the disjunct is infeasible. Otherwise it is the MCKP
\begin{align*}
 L(\theta):=\max_x\quad &\sum_{i,j}v_{ij}x_{ij}\\
 \text{s.t.}\quad &\sum_{i,j}c_{ij}^{\theta}x_{ij}\le C(\theta),\\
 &\sum_jx_{ij}=1,\qquad x\ge0.
\end{align*}

Within a group, plot each option as
$(c_{ij}^{\theta},v_{ij})$. Dominated points do not affect the LP relaxation.
The upper concave hull supplies marginal value-per-cost segments; sorting all
segments by decreasing slope gives the exact greedy solution of the LP
relaxation. At most one group is fractional after merging equal-slope ties in a
consistent way.

Let $\Delta V_i^{\theta}$ be the largest value increase between consecutive
vertices of group $i$'s upper hull, and let
$\Delta V_{\max}^{\theta}:=\max_i\Delta V_i^{\theta}$.

\begin{theorem}[One-item additive certificate]
For a feasible fixed threshold, let $Z^{\rm IP}(\theta)$ and $L(\theta)$ be the
integer and LP optima. Then
\[
  0\le L(\theta)-Z^{\rm IP}(\theta)
  \le \Delta V_{\max}^{\theta}.
\]
A HullRound solution obtained by rounding the sole fractional group toward
the lower-cost adjacent hull vertex is feasible and attains value at least
$L(\theta)-\Delta V_{\max}^{\theta}$.
\end{theorem}

\begin{proof}
The greedy LP fills complete hull segments until the residual capacity cuts at
most one segment. Rounding that segment down cannot increase cost. The value
loss is no larger than the complete value jump of that segment, which is at
most $\Delta V_{\max}^{\theta}$. The integer optimum lies between the rounded
feasible value and the LP upper bound.
\end{proof}

\begin{corollary}[Asymptotic relative gap]
If option values and hull jumps are uniformly bounded while the optimal value
grows linearly in $n$, then the relative fixed-threshold integrality gap is
$O(1/n)$.
\end{corollary}

Upper-hull filtering is safe for the LP calculation but not automatically for
the integer search. A below-hull option can remain part of an optimal integer
combination because capacities are discrete. V3's exact branch-and-bound keeps
all original options except those removed by an integer-safe within-group
dominance rule.

\section{Globally valid exact and anytime search}

For every $\theta\in\B$, denote the fixed-threshold integer optimum by
$Z^{\rm IP}(\theta)$. The robust optimum is
\[
  Z^{\rm robust}=\max_{\theta\in\B}Z^{\rm IP}(\theta).
\]
The exact solver maintains a robust-feasible incumbent value
$Z^{\rm inc}$ and, for each threshold, a valid record $U_{\theta}$ satisfying
$Z^{\rm IP}(\theta)\le U_{\theta}$. Initially, $U_{\theta}=L(\theta)$ whenever
the fixed-threshold LP is feasible and $U_{\theta}=-\infty$ otherwise.
Branch-and-bound may subsequently reduce $U_{\theta}$ or prove the disjunct
optimal.

\begin{proposition}[Global gap invariant]
If every original threshold has a valid upper-bound record, then at every
stage of the search
\[
  Z^{\rm inc}\le Z^{\rm robust}\le
  U^{\rm global}:=\max_{\theta\in\B}U_{\theta}.
\]
Thus a time- or node-limited run reports the valid absolute gap
$U^{\rm global}-Z^{\rm inc}$ and the corresponding scaled relative gap.
\end{proposition}

\begin{proof}
The incumbent is accepted only after direct validation under the original
sorted-$\Gamma$ robust residual, so it is a lower bound. Every robust feasible
selection belongs to at least one original threshold disjunct. The optimum of
that disjunct is bounded by its record, hence the maximum record bounds the
global optimum. Branch-and-bound pruning preserves each record because it
uses only valid LP bounds or proven infeasibility.
\end{proof}

A finite limited-run gap is invalid if some threshold has no upper-bound
record. This accounting rule was a non-negotiable V3 correction: initializing
only the thresholds reached before a time limit can make an apparently small
global gap meaningless.

\section{Exact parametric threshold sweep}

Independent recomputation constructs every array from scratch at every
threshold. V3 also supplies an increasing-$\theta$ sweep. Between consecutive
global breakpoints, each transformed coefficient
$s_{ij}^{\theta}=s_{ij}-\pos{d_{ij}-\theta}$ is affine; at a breakpoint, the
corresponding option changes from the active-deviation expression to its
nominal margin. The sweep updates these coefficients and then recomputes each
group baseline, transformed cost array, and capacity exactly.

Hull reuse is deliberately conservative. A cached group hull is reused only
when the nondominated option set and the complete cost-value point arrays are
unchanged within the stated tolerance. Otherwise it is rebuilt. Optional
validation independently reconstructs sampled states and compares transformed
coefficients, capacities, costs, hulls, and root LP bounds.

\begin{proposition}[Sweep equivalence]
In exact arithmetic, the increasing-threshold sweep produces the same
$s_{ij}^{\theta}$, $c_{ij}^{\theta}$, $C(\theta)$, group hulls, and LP upper
bound as independent fixed-threshold recomputation at every
$\theta\in\B$.
\end{proposition}

\begin{proof}
Each option's update rule is the piecewise-linear identity defining
$s_{ij}^{\theta}$. Baselines, costs, and capacity are then evaluated from their
definitions. A reused hull has an unchanged input point set; a changed point
set is rebuilt. The same deterministic greedy merge therefore receives the
same hull segments as independent recomputation.
\end{proof}

This proposition is an exactness statement, not a universal runtime theorem.
When most group point sets change at most thresholds, hull rebuilding and
fixed-threshold integer search can dominate the saved coefficient work.

\section{Implementation and preserved evidence}

The complete code state is archived at commit
\texttt{2e2829e55243b846e5944965a839eb9b28dd5e9c}. The principal preserved
modules are summarized in \Cref{tab:modules}.

\begin{table}[htbp]
\centering
\caption{V3 code modules preserved in the exact archive.}
\label{tab:modules}
\small
\begin{tabular}{p{0.33\textwidth}p{0.58\textwidth}}
\toprule
Module & Role \\
\midrule
\texttt{certificate.py} & Direct sorted-$\Gamma$ robust-feasibility certificate. \\
\texttt{hull.py}, \texttt{greedy.py} & Group upper hulls and exact fixed-threshold LP relaxation. \\
\texttt{rounding.py}, \texttt{solver.py} & HullRound construction and one-item additive certificate diagnostics. \\
\texttt{exact\_bnb.py} & Fixed-threshold and global full-breakpoint exact branch-and-bound with bound records. \\
\texttt{parametric\_sweep.py} & Exact increasing-threshold state construction, guarded hull reuse, and validation diagnostics. \\
\texttt{local\_budget.py} & Separately scoped segment-local budget extension for small threshold products. \\
\texttt{milp\_baselines.py} & Optional SCIP, HiGHS, Gurobi, and CPLEX baselines with explicit availability reporting. \\
\bottomrule
\end{tabular}
\end{table}

The V3 revision record states that 82 tests passed on 22 May 2026 and that
public and blind manuscript builds passed at that time. It also records that
the complete publication benchmark was not regenerated during the final
revision. The generated CSVs, figures, and manuscript files lived in ignored
local directories and did not survive. Accordingly, this reconstruction does
not repeat unrecoverable V3 runtime tables or claim superiority over mature
MILP solvers. The exact V2/arXiv numerical record remains preserved separately
in the legacy archive.

The surviving design establishes a reproducible route to new evidence:
unit tests exercise certificates and exact-search invariants; smoke drivers
compare independent and swept states; solver scripts report unavailable
backends rather than silently dropping them; and every accepted incumbent is
checked against the original robust constraint.

\section{Scope boundaries and relationship to V4}

V3 is limited to integer $\Gamma$, one global budgeted margin constraint, and
finite exactly-one menus. Fractional uncertainty budgets require an additional
fractional protection term. Multiple independent robust constraints create a
multidimensional threshold product and are not covered by the scalar
decomposition. The pricing application uses own-price demand and does not
establish causal or transaction-level demand effects.

V4 is not a revision of this paper. V3 studies fixed-threshold MCKP hull
geometry, rounding, global enumeration, and exact-search accounting. V4 asks a
different question: how to evaluate a Lagrangian multiplier simultaneously
over the complete threshold family using exactly-one group envelopes. The two
lines share classical robust-MCKP background but have distinct central
theorems, algorithms, and evidence packages.

\section{Conclusion}

The V3 scientific contribution is best understood as a certifying bridge from
finite-menu robust pricing to exact full-breakpoint MCKP decomposition. Its
core guardrails are durable: use every original breakpoint for exactness,
distinguish LP-safe hull filtering from integer-safe pruning, validate every
incumbent in the original model, and report a finite global gap only after
every disjunct has a valid upper-bound record. The parametric sweep preserves
these invariants while avoiding redundant state construction when group point
sets remain unchanged.

This reconstructed PDF restores a readable V3 record without pretending to be
the deleted original. The exact public predecessor, exact V3 code, provenance
notes, checksums, and this reconstruction together form the maintained legacy
archive.

\section*{Archived data and code}

The exact code snapshot, predecessor source, this reconstruction source, and
SHA-256 manifest are stored in the public legacy directory of
\url{https://github.com/eric939/robust_mckp}. The V3 code is independently
anchored by branch \texttt{v3} and tag
\texttt{legacy-v3-code-20260522}.

\begin{thebibliography}{99}

\bibitem[\protect\citeauthoryear{Bertsimas and Sim}{Bertsimas and Sim}{2003}]{BertsimasSim2003}
Bertsimas, D. and Sim, M. (2003).
\newblock Robust discrete optimization and network flows.
\newblock \emph{Mathematical Programming}, 98, 49--71.

\bibitem[\protect\citeauthoryear{Bertsimas and Sim}{Bertsimas and Sim}{2004}]{BertsimasSim2004}
Bertsimas, D. and Sim, M. (2004).
\newblock The price of robustness.
\newblock \emph{Operations Research}, 52, 35--53.

\bibitem[\protect\citeauthoryear{Dyer}{Dyer}{1984}]{Dyer1984}
Dyer, M. E. (1984).
\newblock An $O(n)$ algorithm for the multiple-choice knapsack linear program.
\newblock \emph{Mathematical Programming}, 29, 57--63.

\bibitem[\protect\citeauthoryear{Monaci et al.}{Monaci et al.}{2013}]{Monaci2013}
Monaci, M., Pferschy, U., and Serafini, P. (2013).
\newblock Exact solution of the robust knapsack problem.
\newblock \emph{Computers \& Operations Research}, 40, 2625--2631.

\bibitem[\protect\citeauthoryear{Pisinger}{Pisinger}{1995}]{Pisinger1995}
Pisinger, D. (1995).
\newblock A minimal algorithm for the multiple-choice knapsack problem.
\newblock \emph{European Journal of Operational Research}, 83, 394--410.

\bibitem[\protect\citeauthoryear{Sinha and Zoltners}{Sinha and Zoltners}{1979}]{SinhaZoltners1979}
Sinha, P. and Zoltners, A. A. (1979).
\newblock The multiple-choice knapsack problem.
\newblock \emph{Operations Research}, 27, 503--515.

\end{thebibliography}

\end{document}
